% ==============================================================================
% PRESUPUESTO DESGLOSADO — PROYECTO AUTOVIDA
% CLENCH Software Development · Oferta PRE-2627-01
% Diseño: eje único de importes, Fira Sans, acento magenta reservado al total.
% ==============================================================================
\documentclass[a4paper,10pt]{article}
\usepackage[utf8]{inputenc}
\usepackage[T1]{fontenc}
\usepackage[spanish,es-tabla,es-nodecimaldot]{babel}
\usepackage[left=2.1cm, right=2.1cm, top=3.0cm, bottom=2.6cm,
            headheight=1.6cm, headsep=0.9cm, footskip=1.2cm]{geometry}
\usepackage{graphicx}
\usepackage[table]{xcolor}
\usepackage{tabularx}
\usepackage{longtable}
\usepackage{array}
\usepackage{fancyhdr}
\usepackage{titlesec}
\usepackage{enumitem}
\usepackage{microtype}
\usepackage{lastpage}
\usepackage{tcolorbox}
\tcbuselibrary{skins,breakable}
\usepackage{textcomp}
\usepackage{amsmath}
\usepackage{eso-pic}
\usepackage{letterspace}

% Fira Sans como única familia: humanista, técnica, cifras de altura uniforme.
\usepackage[sfdefault,lf]{FiraSans}
\renewcommand{\familydefault}{\sfdefault}
\linespread{1.12}

% ------------------------------------------------------------------ Paleta
\definecolor{tinta}{HTML}{0E2A33}     % texto y titulares
\definecolor{cian}{HTML}{00A6C7}      % estructura: reglas, numeración
\definecolor{cianoscuro}{HTML}{005D75}
\definecolor{magenta}{HTML}{AD0090}   % reservado al importe total
\definecolor{gris}{HTML}{5A6872}      % texto secundario
\definecolor{banda}{HTML}{EDF5F7}     % cabeceras de tabla
\definecolor{cebra}{HTML}{F7FAFB}     % filas alternas
\definecolor{linea}{HTML}{C9D9DF}
\definecolor{cianclaro}{HTML}{8FC7D6} % texto secundario sobre fondo tinta

\usepackage[hyperfootnotes=false,colorlinks=true,linkcolor=cianoscuro,
            urlcolor=cianoscuro]{hyperref}

% ------------------------------------------------------------------ Rotulos
% Mayusculas con tracking abierto: el recurso tipografico de la marca.
\newcommand{\rotulo}[1]{{\fontsize{8}{10}\selectfont\bfseries\textls[200]{\MakeUppercase{#1}}}}

% Fondo de la portada: tinta a sangre y zocalo cian.
\newcommand{\fondoportada}{%
  \AddToShipoutPictureBG*{%
    \AtPageLowerLeft{\color{tinta}\rule{\paperwidth}{\paperheight}}%
    \AtPageLowerLeft{\color{cian}\rule{\paperwidth}{7mm}}%
  }%
}

% ------------------------------------------------------------------ Titulares
% El número de capítulo es un dato, no un adorno: pequeño, cian, colgado.
\titleformat{\section}[block]
  {\color{tinta}\Large\bfseries}
  {\color{cian}\large\bfseries\thesection\hspace{0.55em}}{0pt}{}
  [\vspace{4pt}\hbox to\linewidth{{\color{cian}\rule[0pt]{2.2cm}{1.4pt}}%
   {\color{linea}\leaders\hrule height 0.5pt\hfill}}]
\titlespacing*{\section}{0pt}{22pt}{4pt}

% El cuadrado cian marca el nivel de subapartado sin gastar un numero.
\titleformat{\subsection}[block]
  {\color{cianoscuro}\normalsize\bfseries}{}{0pt}
  {\raisebox{0.22ex}{\color{cian}\rule{2.8pt}{2.8pt}}\hspace{0.5em}}
\titlespacing*{\subsection}{0pt}{14pt}{4pt}

\titleformat{\subsubsection}[block]
  {\color{cianoscuro}\small\bfseries}{}{0pt}
  {\raisebox{0.2ex}{\color{cian}\rule{2.2pt}{2.2pt}}\hspace{0.45em}}
\titlespacing*{\subsubsection}{0pt}{12pt}{3pt}

\setcounter{secnumdepth}{0}
% Solo los cuatro capitulos de coste directo llevan numero, y es el mismo
% que aparece en sus tablas y en el resumen economico.
\newcommand{\capitulo}[2]{\section{{\color{cian}#1}\hspace{0.55em}#2}}

% ------------------------------------------------------------------ Membrete
\pagestyle{fancy}
\fancyhf{}
\fancyhead[L]{\color{gris}\rotulo{Presupuesto PRE-2627-01\ \ ·\ \ Proyecto AUTOVIDA}}
\fancyhead[R]{\raisebox{-0.35\height}{\includegraphics[height=0.85cm]{images/logo_clench.png}}}
\fancyfoot[L]{\scriptsize\color{gris} CLENCH Software Development\, ·\, E.T.S.I.I.T., c/ Periodista Daniel Saucedo Aranda s/n, 18014 Granada}
\fancyfoot[R]{\scriptsize\color{gris} \thepage\,/\,\pageref{LastPage}}
\renewcommand{\headrulewidth}{0pt}
\renewcommand{\footrulewidth}{0pt}
\fancyheadoffset{0pt}

% Regla fina bajo el membrete, dibujada a mano para controlar su color.
% Regla del membrete en dos tonos: tramo cian de marca y filete fino.
\renewcommand{\headrule}{\vspace{3pt}\hbox to\headwidth{%
  {\color{cian}\rule[0pt]{2.2cm}{1.6pt}}%
  {\color{linea}\leaders\hrule height 0.4pt\hfill}}}
% Filete sobre el pie, para cerrar la caja de texto por abajo.
\renewcommand{\footrule}{\hbox to\headwidth{\color{linea}\leaders\hrule height 0.4pt\hfill}\vspace{4pt}}

% ------------------------------------------------------------------ Tablas
% \hline solo se usa para cerrar bloques de subtotal, de ahi el color y grosor.
\setlength{\arrayrulewidth}{0.7pt}
\arrayrulecolor{cianoscuro}
\renewcommand{\arraystretch}{1.45}
\setlength{\tabcolsep}{7pt}

% Eje único de importes: toda cifra final cae en esta columna.
\newcolumntype{E}{>{\raggedleft\arraybackslash}p{2.75cm}}
\newcolumntype{N}{>{\raggedleft\arraybackslash}p{2.0cm}}
\newcolumntype{H}{>{\raggedleft\arraybackslash}p{1.7cm}}
\newcolumntype{L}{>{\raggedright\arraybackslash}X}

\newcommand{\tablacebra}{\rowcolors{2}{white}{cebra}}
% Cabecera de tabla: banda tinta y texto en blanco, celda a celda para que
% tabularx pueda reprocesar el cuerpo sin romperse.
\newcommand{\cabtxt}[1]{\cellcolor{tinta}\textbf{\color{white}#1}}

% Regla de cierre de bloque, más visible que la de separación de filas.
\newcommand{\reglasubtotal}{\hline}

% Fila de subtotal de capítulo.
\newcommand{\sub}{\rowcolor{white}}

% ------------------------------------------------------------------ Bloques
\newtcolorbox{nota}[1][]{
  colback=cebra, colframe=cebra, coltext=tinta,
  borderline west={2.2pt}{0pt}{cian},
  sharp corners, boxrule=0pt,
  top=3.5mm, bottom=3.5mm, left=4mm, right=4mm,
  fontupper=\small, #1}

\newtcolorbox{condiciones}[1][]{
  colback=cebra, colframe=cebra, coltext=tinta,
  borderline west={2.2pt}{0pt}{cian},
  sharp corners, boxrule=0pt,
  top=3.5mm, bottom=3.5mm, left=4mm, right=4mm,
  fontupper=\small, #1}

\newcommand{\eu}{\,\texteuro}

\begin{document}
\color{tinta}

% ==============================================================================
% PORTADA — tinta a sangre. Identidad primero; la cifra llega en la pagina 2.
% ==============================================================================
\fondoportada
\thispagestyle{empty}
\color{white}

\vspace*{0.4cm}
\noindent\includegraphics[width=5.8cm]{images/logo_clench_portada.png}

\vfill

\noindent{\color{cianclaro}\rotulo{Presupuesto de desarrollo de software}}\\[0.55cm]
\noindent{\color{cian}\rule{3.4cm}{3pt}}\\[0.8cm]
\noindent{\fontsize{52}{54}\selectfont\bfseries AUTOVIDA}

\vfill

\noindent{\color{cian}\rule{\textwidth}{0.5pt}}\\[0.55cm]
\noindent
\begin{minipage}[t]{0.34\textwidth}
  {\color{cian}\rotulo{Destinataria}}\\[0.30cm]
  {\small C.E.E. Fundación Purísima Concepción\\ Granada}
\end{minipage}%
\hfill
\begin{minipage}[t]{0.28\textwidth}
  {\color{cian}\rotulo{Ejecución}}\\[0.30cm]
  {\small 24/09/2026 --- 22/12/2026\\ Fase inicial y tres iteraciones}
\end{minipage}%
\hfill
\begin{minipage}[t]{0.30\textwidth}
  {\color{cian}\rotulo{Oferta}}\\[0.30cm]
  {\small PRE-2627-01, versión 1.0\\ Válida hasta el 11/01/2027}
\end{minipage}

\vspace{1.15cm}
\noindent{\scriptsize\color{cianclaro} CLENCH Software Development\ \ ·\ \ E.T.S.I.I.T., Universidad de Granada}
\vspace*{0.2cm}

\newpage
\color{tinta}

% ==============================================================================
% PÁGINA 2 — MEMBRETE DE LA OFERTA. Abre con la cifra, no con un índice.
% ==============================================================================
\thispagestyle{fancy}

\noindent
\begin{minipage}[t]{0.46\textwidth}
  \vspace{0pt}
  \includegraphics[width=5.2cm]{images/logo_clench.png}\\[0.35cm]
  {\footnotesize\color{gris}
   Desarrollo de software accesible\\
   Granada, España}
\end{minipage}%
\hfill
\begin{minipage}[t]{0.50\textwidth}
  \vspace{0pt}
  {\footnotesize
  \setlength{\tabcolsep}{0pt}
  \begin{tabularx}{\textwidth}{@{}>{\color{gris}}p{3.1cm}X@{}}
  Oferta & \textbf{PRE-2627-01}, versión 1.0\\[2pt]
  Emitida el & 12 de noviembre de 2026\\[2pt]
  Válida hasta el & \textbf{11 de enero de 2027}\\[2pt]
  Destinataria & C.E.E. Fundación Purísima Concepción, Granada\\[2pt]
  Ejecución & 24/09/2026 --- 22/12/2026\\
  \end{tabularx}}
\end{minipage}

\vspace{1.5cm}

\noindent{\color{cian}\rule{\textwidth}{2pt}}\\[0.7cm]

\noindent{\fontsize{30}{34}\selectfont\bfseries Presupuesto de desarrollo}\\[0.45cm]
{\Large\color{cianoscuro} Proyecto AUTOVIDA}

\vspace{1.3cm}

% ---- La cifra. Único uso del magenta en todo el documento. ----
\noindent
\begin{tcolorbox}[colback=banda, colframe=banda, sharp corners, boxrule=0pt,
                  borderline west={4pt}{0pt}{magenta},
                  top=6mm, bottom=6mm, left=6mm, right=6mm]
\begin{minipage}[b]{0.52\textwidth}
  {\footnotesize\color{gris} Importe total, impuestos incluidos}\\[0.25cm]
  {\fontsize{34}{36}\selectfont\bfseries\color{magenta} 35.063,60\eu}
\end{minipage}%
\hfill
\begin{minipage}[b]{0.44\textwidth}
  {\footnotesize\color{gris}
  \setlength{\tabcolsep}{0pt}
  \begin{tabularx}{\textwidth}{@{}Xr@{}}
  Base imponible & \color{tinta}28.978,18\eu\\
  IVA, 21\,\% & \color{tinta}6.085,42\eu\\[3pt]
  \multicolumn{2}{@{}l@{}}{\color{linea}\rule{\textwidth}{0.4pt}}\\[1pt]
  Horas de ingeniería & \color{tinta}1.229,5 h\\
  Coste medio por hora & \color{tinta}23,57\eu\\
  \end{tabularx}}
\end{minipage}
\end{tcolorbox}

\vspace{0.9cm}

\noindent{\small Treinta y cinco mil sesenta y tres euros con sesenta céntimos.}

\vfill

\noindent
\begin{minipage}[t]{0.31\textwidth}
  {\color{cianoscuro}\rotulo{Equipo}}\\[0.25cm]
  {\fontsize{21}{23}\selectfont\bfseries 7}\;{\small\color{gris} personas}\\[0.12cm]
  {\footnotesize\color{gris} Tarifa única para los siete}
\end{minipage}%
\hfill
\begin{minipage}[t]{0.31\textwidth}
  {\color{cianoscuro}\rotulo{Dedicación}}\\[0.25cm]
  {\fontsize{21}{23}\selectfont\bfseries 1.229,5}\;{\small\color{gris} horas}\\[0.12cm]
  {\footnotesize\color{gris} 634,5 ejecutadas, 595 estimadas}
\end{minipage}%
\hfill
\begin{minipage}[t]{0.31\textwidth}
  {\color{cianoscuro}\rotulo{Entregas}}\\[0.25cm]
  {\fontsize{21}{23}\selectfont\bfseries 4}\;{\small\color{gris} hitos}\\[0.12cm]
  {\footnotesize\color{gris} Fase inicial y tres iteraciones}
\end{minipage}

\vspace{1.1cm}

\newpage

\noindent
{\small
\tablacebra
\begin{tabularx}{\textwidth}{@{}LE@{}}
\hline
\cabtxt{Resumen por capítulos} & \cabtxt{Importe}\\
\hline
Personal del equipo, 1.229,5 h a 19,51\eu/h & 23.987,57\eu\\
Inventariable: amortización de equipos informáticos & 164,23\eu\\
Fungible: desplazamientos e impresión & 41,04\eu\\
\hline
\textbf{Costes directos} & \textbf{24.192,84\eu}\\
Costes indirectos de estructura, 13\,\% & 3.145,07\eu\\
\hline
\textbf{Coste de ejecución} & \textbf{27.337,91\eu}\\
Beneficio industrial, 6\,\% & 1.640,27\eu\\
\hline
\textbf{Base imponible} & \textbf{28.978,18\eu}\\
\hline
\end{tabularx}}

% ==============================================================================
\section{Qué se contrata}

AUTOVIDA es una agenda digital accesible para el alumnado del programa PTVAL del C.E.E. Fundación Purísima Concepción, su profesorado y sus familias. Este presupuesto cubre el proyecto completo, desde el análisis de requisitos con el centro hasta la entrega del producto y su transferencia.

\subsection{Clasificación por tipo de gasto}

El presupuesto se presenta de dos formas complementarias: por \textbf{tipo de gasto}, en los capítulos que siguen, y por \textbf{iteraciones}, en el apartado «Reparto por iteraciones». Ambas suman lo mismo.

\vspace{0.35cm}
\noindent{\small
\tablacebra
\begin{tabularx}{\textwidth}{@{}>{\raggedright\arraybackslash}p{2.6cm}L>{\raggedleft\arraybackslash}p{1.5cm}E@{}}
\hline
\cabtxt{Tipo} & \cabtxt{Qué incluye} & \cabtxt{Cap.} & \cabtxt{Importe}\\
\hline
Personal & Las 1.229,5 horas de los siete integrantes, a coste de empresa & 1 & 23.987,57\eu\\
Inventariable & Equipos informáticos, imputados por el desgaste que produce el proyecto y no por su precio de compra & 2 & 164,23\eu\\
Fungible & Desplazamiento de toma de contacto al centro e impresión de entregables (las herramientas software se usan con licencias de estudiante) & 3 & 41,04\eu\\
\hline
\textbf{Costes directos} & & & \textbf{24.192,84\eu}\\
\hline
Indirectos & Estructura de la empresa, compartida con el resto de proyectos (13\,\%) & 4 & 3.145,07\eu\\
Beneficio & Margen empresarial y financiación de inversiones futuras (6\,\%) & 5 & 1.640,27\eu\\
\hline
\textbf{Base imponible} & & & \textbf{28.978,18\eu}\\
\hline
\end{tabularx}}

\vspace{0.4cm}
\noindent La diferencia entre inventariable y fungible no radica en la cuantía, sino en la permanencia del bien: los equipos informáticos continúan en la empresa tras finalizar el proyecto y solo se imputa la fracción de vida útil consumida; los gastos fungibles corresponden a consumos directos de material e impresión y al desplazamiento urbano para las visitas al centro.


\newpage
\subsection{Incluido en el precio}
\begin{itemize}[leftmargin=1.2em, itemsep=2pt, topsep=4pt]
  \item Análisis y especificación de requisitos con participación del centro: personas y escenarios, catálogo de historias de usuario y estimación por \textit{Planning Poker}.
  \item Diseño de la arquitectura ---cliente Flutter, servicio FastAPI, persistencia PostgreSQL--- y prototipado de interfaz en Figma bajo criterios WCAG 2.2 AA y accesibilidad cognitiva.
  \item Implementación en tres iteraciones, cada una con entrega funcional revisable por el centro.
  \item Pruebas unitarias, de integración y de extremo a extremo, más auditorías de accesibilidad con lectores de pantalla y validadores automáticos.
  \item Documentación técnica y de gestión, manual de usuario en lectura fácil y actas de todas las reuniones.
  \item Puesta en marcha del entorno, transferencia y formación básica al personal del centro.
\end{itemize}

\subsection{Fuera del precio}
\begin{itemize}[leftmargin=1.2em, itemsep=2pt, topsep=4pt]
  \item Mantenimiento correctivo y evolutivo posterior a la entrega final, presupuestable aparte.
  \item Explotación en producción: alojamiento más allá del periodo de ejecución, copias de seguridad de larga duración y soporte a personas usuarias.
  \item Adquisición de tabletas o terminales para el alumnado.
  \item Publicación y mantenimiento en tiendas de aplicaciones una vez finalizado el proyecto.
\end{itemize}

% ==============================================================================
\section{El equipo}

Siete personas. Los roles de coordinación del método de trabajo ---Coordinador, Catalogador, Moderador y Presentador--- rotan en cada iteración; los de Gestor de Calidad y Gestor de Usabilidad y Accesibilidad son fijos.

Todo el equipo se factura a una \textbf{tarifa única de 19,51\eu{} por hora}, con independencia del rol que cada persona desempeñe en cada iteración. Los siete integrantes son titulados del mismo nivel, con responsabilidad equivalente sobre el resultado, y los roles rotan por diseño del método de trabajo: diferenciar la tarifa obligaría a reajustarla en cada iteración y haría que el precio dependiese de quién ocupa qué puesto esa semana. Una tarifa única mantiene el precio estable ante cualquier cambio de reparto interno.

\vspace{0.4cm}
\noindent{\small
\tablacebra
\begin{tabularx}{\textwidth}{@{}>{\raggedright\arraybackslash}p{4.4cm}L>{\raggedright\arraybackslash}p{3.5cm}H@{}}
\hline
\cabtxt{Integrante} & \cabtxt{Responsabilidad} & \cabtxt{Categoría} & \cabtxt{\texteuro/h}\\
\hline
José Rodríguez Fernández & Coordinación del proyecto; Presentador en la fase inicial & Área 3, grupo C, nivel 2 & 19,51\\
Julián Carrión Tovar & Usabilidad y accesibilidad (rol fijo); desarrollo frontend & Área 3, grupo C, nivel 2 & 19,51\\
Yeray Rodríguez Navas & Calidad del producto y del proceso (rol fijo) & Área 3, grupo C, nivel 2 & 19,51\\
Pablo de la Torre Roldán & Moderación; desarrollo backend & Área 3, grupo C, nivel 2 & 19,51\\
Pablo Hernández Ibáñez & Comunicación con el centro; desarrollo frontend & Área 3, grupo C, nivel 2 & 19,51\\
Miguel Ángel Luque & Catalogación documental; desarrollo & Área 3, grupo C, nivel 2 & 19,51\\
Amine Azzammouri & Refuerzo de roles rotativos; soporte documental & Área 3, grupo C, nivel 2 & 19,51\\
\hline
\end{tabularx}}

\subsection{De dónde sale la tarifa horaria}

La tarifa horaria representa el coste de empresa por hora de trabajo efectivo de un perfil titulado junior sin experiencia previa, conforme a las tablas salariales del XIX Convenio colectivo estatal de empresas de consultoría y tecnologías de la información (Área 3, grupo C, nivel 2) y las cotizaciones obligatorias a la Seguridad Social (32,15\,\% a cargo de la empresa):

\vspace{0.35cm}
\noindent{\small
\tablacebra
\begin{tabularx}{\textwidth}{@{}L>{\raggedleft\arraybackslash}p{4.2cm}E@{}}
\hline
\cabtxt{Concepto} & \cabtxt{Cálculo / Base} & \cabtxt{Importe}\\
\hline
Retribución bruta anual (convenio) & Salario base (24.461,42\eu) + Plus convenio (2.115,20\eu) & 26.576,62\eu\\
Cotización empresarial Seguridad Social & 32,15\,\% sobre retribución anual & 8.544,38\eu\\
\hline
\textbf{Coste laboral anual para la empresa} & & \textbf{35.121,00\eu}\\
Jornada anual de trabajo efectivo & Según convenio colectivo & 1.800 h\\
\hline
\textbf{Tarifa horaria única} & 35.121,00\eu{} / 1.800 h & \textbf{19,51\eu/h}\\
\hline
\end{tabularx}}

\vspace{0.35cm}
\noindent El tipo de cotización a cargo de la empresa (32,15\,\%) comprende contingencias comunes (23,60\,\%), desempleo (5,50\,\%), contingencias profesionales (1,50\,\%), Mecanismo de Equidad Intergeneracional (0,75\,\%), formación profesional (0,60\,\%) y FOGASA (0,20\,\%).


% ==============================================================================
\capitulo{1}{Personal}

\subsection{Dedicación por integrante y fase, en horas}

\noindent{\small
\tablacebra
\begin{tabularx}{\textwidth}{@{}L H H H H N@{}}
\hline
\cabtxt{Integrante} & \cabtxt{Inicial} & \cabtxt{Iter. 1} & \cabtxt{Iter. 2} & \cabtxt{Iter. 3} & \cabtxt{Total}\\
\hline
José Rodríguez Fernández & 46,5 & 47,5 & 45,5 & 40,5 & 180,0\\
Julián Carrión Tovar & 46,5 & 52,5 & 50,0 & 44,5 & 193,5\\
Yeray Rodríguez Navas & 42,5 & 51,5 & 49,5 & 44,0 & 187,5\\
Pablo de la Torre Roldán & 42,5 & 50,0 & 48,0 & 42,5 & 183,0\\
Pablo Hernández Ibáñez & 42,5 & 45,0 & 43,0 & 38,0 & 168,5\\
Miguel Ángel Luque & 42,5 & 42,5 & 40,5 & 36,0 & 161,5\\
Amine Azzammouri & 42,5 & 40,0 & 38,5 & 34,5 & 155,5\\
\hline
\textbf{Total} & \textbf{305,5} & \textbf{329,0} & \textbf{315,0} & \textbf{280,0} & \textbf{1.229,5}\\
\hline
\end{tabularx}}

\vspace{0.3cm}
\noindent{\footnotesize\color{gris} La dedicación adicional en la fase inicial de José Rodríguez y Julián Carrión responde a las tareas de coordinación, especificación de requisitos y la visita presencial de toma de contacto con el centro.}

\subsection{Importe por integrante}

Con tarifa única, lo que cada persona factura depende solo de las horas que dedica.

\noindent{\small
\tablacebra
\begin{tabularx}{\textwidth}{@{}L N N E@{}}
\hline
\cabtxt{Integrante} & \cabtxt{Horas} & \cabtxt{\texteuro/h} & \cabtxt{Importe}\\
\hline
José Rodríguez Fernández & 180,0 & 19,51 & 3.511,80\eu\\
Julián Carrión Tovar & 193,5 & 19,51 & 3.775,19\eu\\
Yeray Rodríguez Navas & 187,5 & 19,51 & 3.658,13\eu\\
Pablo de la Torre Roldán & 183,0 & 19,51 & 3.570,33\eu\\
Pablo Hernández Ibáñez & 168,5 & 19,51 & 3.287,44\eu\\
Miguel Ángel Luque & 161,5 & 19,51 & 3.150,87\eu\\
Amine Azzammouri & 155,5 & 19,51 & 3.033,81\eu\\
\hline
\textbf{Capítulo 1} & \textbf{1.229,5} & \textbf{19,51} & \textbf{23.987,57\eu}\\
\hline
\end{tabularx}}

\subsection{En qué se emplean esas horas}

\noindent{\small
\tablacebra
\begin{tabularx}{\textwidth}{@{}L N E@{}}
\hline
\cabtxt{Actividad} & \cabtxt{\%} & \cabtxt{Horas}\\
\hline
Implementación \textit{frontend} y \textit{backend} & 32,0\,\% & 393,5\\
Reuniones de coordinación, seguimiento y cliente & 27,0\,\% & 332,0\\
Pruebas y auditorías de accesibilidad & 12,0\,\% & 147,5\\
Análisis, requisitos y gestión del \textit{backlog} & 10,0\,\% & 123,0\\
Documentación, actas y entregables & 10,0\,\% & 123,0\\
Diseño de arquitectura y prototipado de interfaz & 9,0\,\% & 110,5\\
\hline
\textbf{Total} & \textbf{100,0\,\%} & \textbf{1.229,5}\\
\hline
\end{tabularx}}

\vspace{0.4cm}
\noindent Las 332 horas de reuniones corresponden a las once sesiones documentadas en actas entre el 1 de octubre y el 11 de noviembre de 2026 ---45 horas por persona--- más las sesiones previstas para las iteraciones 2 y 3.

% ==============================================================================
\capitulo{2}{Material inventariable: equipos informáticos}

Los equipos son material \textbf{inventariable}: bienes duraderos que siguen siendo de la empresa cuando el proyecto termina. Por eso no se imputa su precio de compra, sino la parte de su valor que el proyecto consume. El prorrateo combina tres datos:

\begin{enumerate}[leftmargin=1.4em, itemsep=3pt, topsep=4pt]
  \item El \textbf{coste real del equipo hoy}, tomado de precios de venta al público vigentes en 2026 para equipos equivalentes.
  \item Su \textbf{duración estimada}: cuatro años, que es el periodo más corto admitido por la tabla de amortización del impuesto sobre sociedades para equipos de procesos de información, cuyo coeficiente lineal máximo es del 25\,\% anual.
  \item El \textbf{tiempo de uso en el proyecto}, en horas, no en meses de calendario: un portátil se imputa por las horas de trabajo efectivo y los dispositivos de prueba por las horas de pruebas.
\end{enumerate}

\vspace{0.35cm}
\begin{center}
{\color{cianoscuro}
 \large Importe imputado\ \ =\ \
 \begin{tabular}[c]{@{}c@{}}
 coste actual del equipo\\[1pt]
 \hline
 \rule{0pt}{14pt}4 años $\times$ 1.800 h
 \end{tabular}
 \ \ $\times$\ \ horas de uso}
\end{center}
\vspace{0.35cm}

\noindent{\small
\tablacebra
\begin{tabularx}{\textwidth}{@{}L>{\raggedleft\arraybackslash}p{1.8cm}>{\raggedleft\arraybackslash}p{1.7cm}>{\raggedleft\arraybackslash}p{1.7cm}E@{}}
\hline
\cabtxt{Equipo} & \cabtxt{Coste hoy} & \cabtxt{\texteuro/h} & \cabtxt{Uso} & \cabtxt{Importe}\\
\hline
Portátil de desarrollo, 16 GB de RAM y SSD de 512 GB (7 ud.) & 755,16\eu & 0,1049 & 1.229,5 h & 128,95\eu\\
Monitor externo de 24\,pulgadas (7 ud.) & 129,00\eu & 0,0179 & 1.229,5 h & 22,03\eu\\
Tableta Android de 10\,pulgadas para pruebas en el dispositivo objetivo (2 ud.) & 199,00\eu & 0,0276 & 147,5 h & 8,15\eu\\
Terminal Android para validación con TalkBack (1 ud.) & 249,00\eu & 0,0346 & 147,5 h & 5,10\eu\\
\hline
\textbf{Capítulo 2} & & & & \textbf{164,23\eu}\\
\hline
\end{tabularx}}

\vspace{0.4cm}
\noindent{\footnotesize\color{gris} Los portátiles y los monitores se imputan por las 1.229,5 horas de trabajo del equipo, porque cada hora de trabajo ocupa un puesto completo. Las tabletas y el terminal de pruebas se imputan por las 147,5 horas dedicadas a pruebas y auditorías de accesibilidad, que es el único momento en que se usan, y por las unidades empleadas.}




% ==============================================================================
\capitulo{3}{Material fungible: desplazamientos e impresión}

Gastos corrientes directamente necesarios para la ejecución del proyecto que se agotan durante el mismo.

Las herramientas y plataformas software empleadas para la coordinación, control de versiones, diseño y gestión del proyecto (Jira, GitHub, Figma, etc.) se cubren mediante licencias educativas y planes para estudiantes sin coste económico repercutible al proyecto. Asimismo, las tecnologías base de desarrollo (Flutter, FastAPI, PostgreSQL) y los recursos gráficos de ARASAAC (licencia CC BY-NC-SA) corresponden a software libre y licencias abiertas.

Por consiguiente, esta partida recoge exclusivamente el desplazamiento urbano necesario para la visita presencial al centro educativo y los gastos de reprografía y encuadernación de las entregas documentales:

\vspace{0.4cm}
\noindent{\small
\tablacebra
\begin{tabularx}{\textwidth}{@{}L>{\raggedright\arraybackslash}p{5.5cm}E@{}}
\hline
\cabtxt{Concepto} & \cabtxt{Cálculo} & \cabtxt{Importe}\\
\hline
Desplazamiento al centro (visita presencial) & 1 persona × 1 visita (4 km ida y vuelta) × 0,26\eu/km & 1,04\eu\\
Impresión y encuadernación de entregables & Memorias de seguimiento y documentación final & 40,00\eu\\
\hline
\textbf{Capítulo 3} & & \textbf{41,04\eu}\\
\hline
\end{tabularx}}

% ==============================================================================
\section{Estructura y beneficio}

Los importes añadidos a los costes directos corresponden a los gastos generales de soporte de la empresa y al margen empresarial de la actividad.

\subsection{Costes indirectos de estructura}

Son los gastos que la empresa soporta para hacer viable el proyecto pero que \textbf{comparte con el resto de su actividad}: instalaciones, suministros, asesoría y administración, seguros de responsabilidad civil y plataformas comunes. Al no ser imputables a tareas concretas, se repercuten mediante un 13\,\% sobre los costes directos, coeficiente habitual de referencia en contratos de servicios de ingeniería y consultoría. Su desglose orientativo es:

\vspace{0.4cm}
\noindent{\small
\tablacebra
\begin{tabularx}{\textwidth}{@{}L>{\raggedright\arraybackslash}p{4.4cm}E@{}}
\hline
\cabtxt{Concepto} & \cabtxt{Cálculo} & \cabtxt{Importe}\\
\hline
Espacio de trabajo, suministros y limpieza & 5,0\,\% de los directos & 1.209,64\eu\\
Administración, contabilidad y facturación & 3,0\,\% & 725,79\eu\\
Formación interna y actualización técnica & 2,0\,\% & 483,86\eu\\
Seguro de responsabilidad civil profesional & 1,5\,\% & 362,89\eu\\
Sistemas corporativos y gestión documental & 1,5\,\% & 362,89\eu\\
\hline
\textbf{Capítulo 4} & \textbf{13,0\,\%} & \textbf{3.145,07\eu}\\
\hline
\end{tabularx}}

\subsection{Beneficio industrial}

Se aplica un 6\,\% sobre el coste de ejecución (27.337,91\eu), tipo de referencia estándar en la contratación del sector público español. Retribuye la actividad empresarial a precio cerrado, compensando posibles oscilaciones operativas de esfuerzo interno y posibilitando la reinversión y sostenimiento tecnológico del equipo.

\vspace{0.4cm}
\noindent{\small
\tablacebra
\begin{tabularx}{\textwidth}{@{}L>{\raggedright\arraybackslash}p{4.4cm}E@{}}
\hline
\cabtxt{Concepto} & \cabtxt{Cálculo} & \cabtxt{Importe}\\
\hline
Beneficio industrial & 6,0\,\% sobre coste de ejecución & 1.640,27\eu\\
\hline
\textbf{Capítulo 5} & & \textbf{1.640,27\eu}\\
\hline
\end{tabularx}}

\subsection{Resumen económico}

\vspace{0.2cm}
\noindent{\small
\tablacebra
\begin{tabularx}{\textwidth}{@{}>{\raggedright\arraybackslash}p{1.0cm}L>{\raggedleft\arraybackslash}p{2.0cm}E@{}}
\hline
\cabtxt{Cap.} & \cabtxt{Concepto} & \cabtxt{\% s/ base} & \cabtxt{Importe}\\
\hline
1 & Personal, 1.229,5 h a 19,51\eu/h & 82,78\,\% & 23.987,57\eu\\
2 & Inventariable: equipos informáticos & 0,57\,\% & 164,23\eu\\
3 & Fungible: desplazamientos e impresión & 0,14\,\% & 41,04\eu\\
\hline
 & \textbf{Costes directos} & \textbf{83,49\,\%} & \textbf{24.192,84\eu}\\
4 & Costes indirectos de estructura, 13\,\% & 10,85\,\% & 3.145,07\eu\\
\hline
 & \textbf{Coste de ejecución} & \textbf{94,34\,\%} & \textbf{27.337,91\eu}\\
5 & Beneficio industrial, 6\,\% & 5,66\,\% & 1.640,27\eu\\
\hline
 & \textbf{Base imponible} & \textbf{100,00\,\%} & \textbf{28.978,18\eu}\\
 & IVA, 21\,\% & & 6.085,42\eu\\
\hline
 & \textbf{Importe total} & & \textbf{35.063,60\eu}\\
\hline
\end{tabularx}}

% ==============================================================================
\section{Reparto por iteraciones}

El proyecto se organiza en una fase inicial de análisis y tres iteraciones de desarrollo, cada una cerrada con una entrega revisable por el centro. Los importes de cada fase reparten los capítulos 2 a 5 en proporción a las horas imputadas.

\subsection{Calendario}

\noindent{\small
\tablacebra
\begin{tabularx}{\textwidth}{@{}>{\raggedright\arraybackslash}p{2.4cm}>{\raggedright\arraybackslash}p{3.1cm}L>{\raggedright\arraybackslash}p{1.9cm}@{}}
\hline
\cabtxt{Fase} & \cabtxt{Periodo} & \cabtxt{Resultado} & \cabtxt{Estado}\\
\hline
Fase inicial & 24/09 --- 22/10/2026 & Requisitos, personas y escenarios, catálogo de historias de usuario estimadas, manual de coordinación y propuesta técnica & Ejecutada\\
Iteración 1 & 23/10 --- 11/11/2026 & Arquitectura, API REST y cliente Flutter cubriendo 32 historias de usuario; primera auditoría de accesibilidad & Ejecutada\\
Iteración 2 & 12/11 --- 03/12/2026 & Agenda completa, pasos guiados de actividad, ajustes de accesibilidad por perfil y validación con usuarios & Estimada\\
Iteración 3 & 04/12 --- 22/12/2026 & Funcionalidades alternativas, corrección de la auditoría final, manual en lectura fácil y transferencia & Estimada\\
\hline
\end{tabularx}}

\subsection{Importe por fase}

\noindent{\small
\tablacebra
\begin{tabularx}{\textwidth}{@{}L H>{\raggedleft\arraybackslash}p{2.2cm}>{\raggedleft\arraybackslash}p{2.0cm}>{\raggedleft\arraybackslash}p{2.0cm}E@{}}
\hline
\cabtxt{Fase} & \cabtxt{Horas} & \cabtxt{Personal} & \cabtxt{Invent. y fung.} & \cabtxt{Estruct.} & \cabtxt{Base imp.}\\
\hline
Fase inicial & 305,5 & 5.960,31\eu & 51,00\eu & 781,47\eu & 7.200,35\eu\\
Iteración 1 & 329,0 & 6.418,79\eu & 54,93\eu & 841,58\eu & 7.754,22\eu\\
Iteración 2 & 315,0 & 6.145,65\eu & 52,59\eu & 805,77\eu & 7.424,25\eu\\
Iteración 3 & 280,0 & 5.462,82\eu & 46,75\eu & 716,25\eu & 6.599,36\eu\\
\hline
\textbf{Total} & \textbf{1.229,5} & \textbf{23.987,57\eu} & \textbf{205,27\eu} & \textbf{3.145,07\eu} & \textbf{28.978,18\eu}\\
\hline
\end{tabularx}}

\vspace{0.3cm}
\noindent{\footnotesize\color{gris} La cuarta columna recoge la amortización de inventariable y los gastos fungibles imputables a la fase. La base imponible de cada fase incluye su parte proporcional del beneficio industrial del 6\,\%, que suma 1.640,27\eu{} en el conjunto del proyecto.}

\subsection{Cómo se han estimado las iteraciones 2 y 3}

\begin{itemize}[leftmargin=1.2em, itemsep=3pt, topsep=4pt]
  \item \textbf{Iteración 2, 315 horas.} Estimada por analogía con la iteración 1, que consumió 329 horas para cubrir 32 historias de usuario. El alcance previsto es comparable en volumen, pero el equipo parte de una arquitectura, un entorno y unas convenciones ya establecidos, lo que reduce la dedicación en torno a un 4\,\%.
  \item \textbf{Iteración 3, 280 horas.} Un 15\,\% por debajo de la iteración 1. Las historias pendientes son funcionalidades alternativas de menor complejidad, aunque se reserva dedicación para corregir la auditoría final de accesibilidad, redactar el manual en lectura fácil y transferir el producto al centro.
  \item Ambas estimaciones mantienen el reparto relativo de horas entre integrantes observado en la iteración 1 y asumen la continuidad del equipo y de los roles.
\end{itemize}

% ==============================================================================

\newpage
\section{Calendario de facturación}

Se factura al cierre y aceptación de cada fase, contra la entrega revisable correspondiente.

\vspace{0.4cm}
\noindent{\small
\tablacebra
\begin{tabularx}{\textwidth}{@{}L>{\raggedright\arraybackslash}p{2.2cm}N N E@{}}
\hline
\cabtxt{Hito} & \cabtxt{Fecha} & \cabtxt{Base} & \cabtxt{IVA 21\,\%} & \cabtxt{Total}\\
\hline
Cierre de la fase inicial & 22/10/2026 & 7.200,35\eu & 1.512,07\eu & 8.712,42\eu\\
Entrega de la iteración 1 & 11/11/2026 & 7.754,22\eu & 1.628,39\eu & 9.382,61\eu\\
Entrega de la iteración 2 & 03/12/2026 & 7.424,25\eu & 1.559,09\eu & 8.983,34\eu\\
Entrega final y transferencia & 22/12/2026 & 6.599,36\eu & 1.385,87\eu & 7.985,23\eu\\
\hline
\textbf{Total} & & \textbf{28.978,18\eu} & \textbf{6.085,42\eu} & \textbf{35.063,60\eu}\\
\hline
\end{tabularx}}
% ==============================================================================
\vspace{0.8cm}
\section{Fuentes}

Las cifras normativas y de referencia aplicadas en esta oferta proceden de fuentes oficiales y públicas:

\subsection{Coste de personal}
\begin{itemize}[leftmargin=1.2em, itemsep=3pt, topsep=3pt]
  \item \textbf{Convenio colectivo y tablas salariales (2026):} XIX Convenio colectivo estatal de empresas de consultoría y tecnologías de la información (BOE núm. 100, de 24 de abril de 2026; BOE-A-2026-9024). Área 3, grupo C, nivel 2 (titulados junior): salario base anual de 24.461,42\eu{} y plus de convenio de 2.115,20\eu{} (total 26.576,62\eu). Jornada máxima de 1.800 horas de trabajo efectivo en cómputo anual (BOE-A-2025-7766).
  \item \textbf{Cotización a la Seguridad Social (2026):} Orden PJC/297/2026, con tipo global empresarial del 32,15\,\% para contrato indefinido (contingencias comunes 23,60\,\%, desempleo 5,50\,\%, accidentes de trabajo y enfermedad profesional 1,50\,\%, MEI 0,75\,\%, formación profesional 0,60\,\% y FOGASA 0,20\,\%).
\end{itemize}

\subsection{Material inventariable y fungible}
\begin{itemize}[leftmargin=1.2em, itemsep=3pt, topsep=3pt]
  \item \textbf{Amortización de equipos:} Tabla de amortización simplificada (art. 12.1.a de la Ley 27/2014 del Impuesto sobre Sociedades): coeficiente lineal máximo del 25\,\% anual (vida útil de 4 años para equipos informáticos), imputados exclusivamente por las horas de uso efectivo en el proyecto.
  \item \textbf{Herramientas y licencias:} Software de gestión y desarrollo cubierto íntegramente mediante licencias educativas y de estudiante sin coste para el cliente.
  \item \textbf{Kilometraje:} 0,26\eu/km según la cuantía exenta de IRPF para locomoción (Orden HFP/792/2023).
\end{itemize}

\subsection{Porcentajes aplicados}
\begin{itemize}[leftmargin=1.2em, itemsep=3pt, topsep=3pt]
  \item \textbf{Costes indirectos de estructura (13\,\%):} Coeficiente de referencia habitual para gastos generales de estructura en contratos de servicios y proyectos cofinanciados.
  \item \textbf{Beneficio industrial (6\,\%):} Margen empresarial estándar de referencia en la contratación del sector público español.
\end{itemize}


\newpage
\section{Condiciones}

\begin{enumerate}[leftmargin=1.4em, itemsep=4pt, topsep=4pt]
  \item \textbf{Validez.} Esta oferta es válida durante 60 días naturales desde su emisión, hasta el 11 de enero de 2027 inclusive. Pasado ese plazo sin aceptación expresa, los importes se revisarán conforme a las tarifas vigentes.
  \item \textbf{Impuestos.} Los importes de los capítulos se expresan sin IVA. El impuesto se repercute al tipo general del 21\,\% vigente a la fecha de emisión; cualquier modificación legal del tipo se trasladará a la facturación pendiente.
  \item \textbf{Forma de pago.} Transferencia bancaria a 30 días desde la emisión de cada factura, según el calendario de hitos.
  \item \textbf{Revisión de precios.} No se aplica revisión dentro del periodo de ejecución previsto.
  \item \textbf{Cambios de alcance.} Toda petición que amplíe el alcance descrito se tramita por el procedimiento de control de cambios del manual de coordinación, con estimación previa de horas e importe y aprobación escrita del centro antes de ejecutarse.
  \item \textbf{Propiedad intelectual.} El código fuente y la documentación se entregan al centro con licencia de uso perpetua, no exclusiva y sin coste adicional. Los pictogramas de ARASAAC conservan su licencia Creative Commons original.
  \item \textbf{Protección de datos.} El tratamiento de datos del alumnado se ajusta al RGPD y a la LOPDGDD. No se incluye la designación de delegado de protección de datos, que corresponde al centro como responsable del tratamiento.
  \item \textbf{Garantía.} Treinta días naturales desde la entrega final para la corrección de defectos imputables al desarrollo, sin coste adicional.
\end{enumerate}

\vfill

\noindent
\begin{minipage}[t]{0.47\textwidth}
{\footnotesize\color{gris} Por CLENCH Software Development}\\[1.8cm]
{\color{linea}\rule{\textwidth}{0.4pt}}\\[2pt]
{\small José Rodríguez Fernández}\\
{\footnotesize\color{gris} Coordinador del proyecto\\ Granada, 12 de noviembre de 2026}
\end{minipage}%
\hfill
\begin{minipage}[t]{0.47\textwidth}
{\footnotesize\color{gris} Conforme, por el C.E.E. Fundación Purísima Concepción}\\[1.8cm]
{\color{linea}\rule{\textwidth}{0.4pt}}\\[2pt]
{\small Dirección del centro}\\
{\footnotesize\color{gris} Nombre y fecha}
\end{minipage}

\end{document}
